\documentclass{beamer}
\usepackage[dvipsnames]{xcolor}
\usetheme[hideothersubsections, width=1.7cm]{Hannover}
\usecolortheme{spruce}

\usepackage{graphicx}
\usepackage{animate}
\usepackage{bm}
\usepackage{graphicx}

% Datos
\title{Dinámica Peatonal: Caminata al Azar}
\author{Grupo 5}
\institute{ITBA}
\date{} % sin fecha

% ---- Colores de prueba (bien contrastantes) ----
% Color base para bullets, numeración, títulos por defecto
\setbeamercolor{structure}{fg=blue}

% Sidebar: secciones
\usebeamercolor{title in sidebar}
\colorlet{SidebarTitleFG}{title in sidebar.fg}
%\setbeamercolor{section in sidebar}{fg=white, bg=red}          % sección activa -
\setbeamercolor{section in sidebar shaded}{fg=SidebarTitleFG!70!black} % secciones inactivas
\setbeamercolor{title in sidebar}{fg=SidebarTitleFG!70!black}
\setbeamercolor{author in sidebar}{fg=SidebarTitleFG!90!black}

% Sidebar: subsecciones
\usebeamercolor{subsection in sidebar}
\colorlet{SidebarSubsectionFG}{subsection in sidebar.fg}
\setbeamercolor{subsection in sidebar}{fg=white}           % subsección activa
\setbeamercolor{subsection in sidebar shaded}{}    % subsecciones inactivas

% Topbar (headline con palettes)

\begin{document}

% Portada
  \begin{frame}
    \titlepage
    \begin{center}
      \small INTEGRANTES: MARTINA SCHVARTZ TALLONE, PATRICK LUCA TORLASCHI y SERGIO SMIRNOFF
    \end{center}
  \end{frame}

% ------------------ Introducción ------------------
  \section{Introducción}
  \begin{frame}
    \centering
    \colorbox{green!20!gray!30}{%
      \begin{minipage}{0.85\textwidth}
        \centering
        \vspace{0.5cm}
        {\Large\ Introducción}
        \vspace{0.5cm}
      \end{minipage}
    }
  \end{frame}

  \subsection{Modelo Matemático}
  \begin{frame}{Social Force Model}
    \begin{itemize}
      \item Ecuación de Movimiento:
      \vspace{0.2cm}
      \\\( m_i\frac{d\mathbf{v}_i}{dt}=\mathbf{f}_{ap}+\sum_{j\ne i}{\mathbf{f}_{ij}}\)
      \vspace{0.5cm}
      \item Fuerza de Interacción:
      \vspace{0.1cm}
      \begin{columns}
        \begin{column}{0.6\textwidth}
            \\\(\mathbf{f}_{ij}=\textit{k}_ng(r_{ij}-d_{ij})\mathbf{n}_{ij}+\textit{k}_tg(r_{ij}-d_{ij})\Delta \textit{v}^t_{ji}\mathbf{t}_{ij}\)
            \vspace{0.2cm}
            \\donde \( g(x)=\begin{cases}
                              0 & d_{ij} \ge r_{ij} \\
                              x & d_{ij} < r_{ij}
              \end{cases}\)
        \end{column}
        \begin{column}{0.4\textwidth}
          \includegraphics[width=0.3\linewidth]{imagen/interaction.png}
        \end{column}
      \end{columns}
      \vspace{0.5cm}
      \item Fuerza Autopropulsora:
      \vspace{0.2cm}
      \\\(\mathbf{f}_{ap} = m_i\frac{v^0_i\mathbf{e}^0_i(t)-\mathbf{v}_i(t)}{\tau_i}\)
      \vspace{0.2cm}
      \\donde \(v^0\) es la velocidad deseada y \(\tau_i\) el tiempo característico
    \end{itemize}
  \end{frame}

% ------------------ Implementación ------------------

  \section{Implementación}
  \begin{frame}
    \centering
    \colorbox{green!20!gray!30}{%
      \begin{minipage}{0.85\textwidth}
        \centering
        \vspace{0.5cm}
        {\Large\ Implementación}
        \vspace{0.5cm}
      \end{minipage}
    }
  \end{frame}

  \subsection{Implementación}
  \begin{frame}{UML}
    \begin{columns}
      \begin{column}{0.4\textwidth}
        Clases:
        \begin{itemize}
          \item Peaton
          \item Simulation
          \item CellIndexMethod
          \item Integrator: Beeman
        \end{itemize}
      \end{column}

      \begin{column}{0.6\textwidth}
        \begin{figure}
          \centering
          \includegraphics[width=0.95\linewidth]{imagen/uml.jpeg}
          \label{fig:placeholder}
        \end{figure}
        \vspace{0.2cm}
      \end{column}
    \end{columns}
  \end{frame}


% ------------------ Simulaciones ------------------

  \section{Simulaciones}
  \begin{frame}
    \centering
    \colorbox{green!20!gray!30}{%
      \begin{minipage}{0.85\textwidth}
        \centering
        \vspace{0.5cm}
        {\Large\ Simulaciones}
        \vspace{0.5cm}
      \end{minipage}
    }
  \end{frame}

  \subsection{Geometría del Sistema}
  \begin{frame}{Geometría del Sistema}

    \begin{columns}
      \begin{column}{0.4\textwidth}
        %FOTO SISTEMA (ANIMACIÓN)
        \begin{itemize}
          \item \(L=6.0m\)
          \vspace{0.2cm}
          \item \(v^0=1.7\frac{m}{s}\)
          \vspace{0.2cm}
          \item \(dt=0.0001\ s\)
          \vspace{0.2cm}
          \item t característico \(=0.5\ s\)
          \vspace{0.2cm}
          \item \(\textit{k}_n=1.2\cdot10^5\frac{kg}{s^2}\)
          \vspace{0.2cm}
          \item \(\textit{k}_t=2.4\cdot10^5\frac{kg}{ms}\)
        \end{itemize}
        \vspace{0.5cm}
        \medium {Input:}
        \begin{itemize}
          \item \(\phi=\frac{\sum^N_i{\pi \cdot r^2_i}}{L^2}\), $N\in[10,220]$
        \end{itemize}

      \end{column}

      \begin{column}{0.6\textwidth}
        \begin{figure}
          \centering
          \includegraphics[width=1\linewidth]{imagen/image.png}
          \label{fig:placeholder}
        \end{figure}
      \end{column}
    \end{columns}

  \end{frame}

  \subsection{Definición de Observables}
  \begin{frame}{Definición de Observables}
    \begin{itemize}
      \item \(Q_{\text{prom}} = \langle Q \rangle_{10\,\text{corridas}} \pm \sigma(Q)\)
      \vspace{0.2cm}
      \\\(Q=\frac{dN(t)}{dt}\)
      \vspace{0.2cm}
      \\ \(N(t)=cant. contactos\ únicos\)
      % ACLARAR QUE Q ES EL PROMEDIO DE 10 CORRIDAS
      \vspace{0.5cm}
      \item \(\text{CDSF}(\bm{\tau}) = P(T \ge \bm{\tau}) = 1 - F(\bm{\tau})\)
      \vspace{0.1cm}
      \\ \(\bm{\tau}_i=t_{i+1}-t_i\)
      % NUEVO ITEM AJUSTE DE ALPHA EXPLICACIÓN BREVE
      \vspace{0.2cm}
      \item $ \alpha $ y $ \lambda $: parámetros Ley de Potencias y Exponencial respectivamente.
    \end{itemize}
  \end{frame}

  % ------------------ Resultados Base ------------------
  \section{Resultados}
  
  \subsection{Resultados}
  \begin{frame}
    \centering
    \colorbox{green!20!gray!30}{%
      \begin{minipage}{0.85\textwidth}
        \centering
        \vspace{0.5cm}
        {\Large\ Resultados}
        \vspace{0.5cm}
      \end{minipage}
    }
  \end{frame}

  \subsubsection{Animaciones}
  \begin{frame}{Animaciones}

  \end{frame}
  \begin{frame}{Animaciones}

    \begin{columns}
      \begin{column}{0.5\textwidth}
        \medium {N = 72}
        \includegraphics[width=1\linewidth]{imagen/ani_N72_L6.png}
        \url {https://www.youtube.com/watch?v=2i7REk3KfTc&t=36s}
      \end{column}

      \begin{column}{0.5\textwidth}
        \medium {N = 220}
        \includegraphics[width=1\linewidth]{imagen/ani_N220_L6.png}
        \url {https://www.youtube.com/watch?v=yL0JS209mKc&t=66s}
      \end{column}
    \end{columns}

  \end{frame}

  \subsubsection{Contactos Únicos con Agente Fijo}
  \begin{frame}{Contactos Únicos con Agente Fijo}
    \begin{figure}
      \centering
      \includegraphics[width=1\linewidth]{imagen/accumulated_contacts_N_9_72_219.png}
      \label{fig:placeholder}
    \end{figure}
  \end{frame}

  \subsubsection{Scanning Rate}
  \begin{frame}{Scanning Rate}
    \begin{figure}
      \centering
      \includegraphics[width=1\linewidth]{imagen/scanning_rate_vs_phi.png}
      \label{fig:placeholder}
    \end{figure}
  \end{frame}

  \subsubsection{Aproximaciones}

  \begin{frame}{Aproximaciones}
    \begin{columns}
      \begin{column}{0.5\textwidth}
        \centering
        \medium{Aproximación Power Law}
        \begin{figure}
          \small{N=39}
          \includegraphics[width=\linewidth]{imagen/distribution_N39.png}
          \label{fig:dist_N39}
        \end{figure}
      \end{column}
      \begin{column}{0.5\textwidth}
        \centering
        \medium{Aproximación Exponencial}
        \begin{figure}
          \small{N=54}
          \includegraphics[width=\linewidth]{imagen/distribution_N54.png}
          \label{fig:dist_N54}
        \end{figure}
      \end{column}
    \end{columns}
  \end{frame}


  \subsubsection{Exponentes Aproximaciones}
  \begin{frame}{Exponentes: Ley de Potencias y Exponencial}
    \begin{columns}
      \begin{column}{0.5\textwidth}
        \centering
        \medium{Exponente de la Ley de Potencias}
        \vspace{0.2cm}
        \begin{figure}
          \centering
          \includegraphics[width=\linewidth]{imagen/alpha_vs_phi.png}
          \label{fig:alpha_vs_phi}
        \end{figure}  
      \end{column}
      \begin{column}{0.5\textwidth}
        \centering
        \medium{Exponente de Exponencial}
        \vspace{0.2cm}
        \begin{figure}
          \centering
          \includegraphics[width=\linewidth]{imagen/lambda_vs_phi.png}
          \label{fig:lambda_vs_phi}
        \end{figure}
      \end{column}
    \end{columns}
  \end{frame}


  % ------------------ Resultados V = 2.5 ------------------

  \subsection{Resultados V 2.5}
  \begin{frame}
    \centering
    \colorbox{green!20!gray!30}{%
      \begin{minipage}{0.85\textwidth}
        \centering
        \vspace{0.5cm}
        {\Large\ Resultados V 2.5}
        \vspace{0.5cm}
      \end{minipage}
    }
  \end{frame}

  \subsubsection{Animaciones}
  \begin{frame}{Animaciones}
  \end{frame}
  \begin{frame}{Animaciones}
    \begin{columns}
      \begin{column}{0.5\textwidth}
        \medium {N = 69}
        \includegraphics[width=1\linewidth]{imagenv25/ani_N69_L6.png}
        \url {https://www.youtube.com/watch?v=zeKmTOo-9jc}
      \end{column}
      \begin{column}{0.5\textwidth}
        \medium {N = 219}
        \includegraphics[width=1\linewidth]{imagenv25/ani_N219_L6.png}
        \url {https://www.youtube.com/watch?v=DpU_IK-XHEc}
      \end{column}
    \end{columns}
  \end{frame}


  \subsubsection{Contactos Únicos con Agente Fijo}
  \begin{frame}{Contactos Únicos con Agente Fijo}
    \begin{figure}
      \centering
      \includegraphics[width=1\linewidth]{imagenv25/accumulated_contacts_N_9_69_219.png}
      \label{fig:placeholder}
    \end{figure}
  \end{frame}

  \subsubsection{Scanning Rate}
  \begin{frame}{Scanning Rate}
    \begin{figure}
      \centering
      \includegraphics[width=1\linewidth]{imagenv25/scanning_rate_vs_phi.png}
      \label{fig:placeholder}
    \end{figure}
  \end{frame}

  \subsubsection{Aproximaciones}

  \begin{frame}{Aproximaciones}
    \begin{columns}
      \begin{column}{0.5\textwidth}
        \medium{Aproximación Power Law}
        \begin{figure}
          \centering
          \small{N=39}
          \includegraphics[width=1\linewidth]{imagenv25/distribution_N39.png}
          \label{fig:dist_N39}
        \end{figure}
      \end{column}
      \begin{column}{0.5\textwidth}
        \medium{Aproximación Exponencial}
        \begin{figure}
          \centering
          \small{N=219}
          \includegraphics[width=1\linewidth]{imagenv25/distribution_N219.png}
          \label{fig:dist_N219}
        \end{figure}
      \end{column}
    \end{columns}
  \end{frame}


  \subsubsection{Exponentes Aproximaciones}
  \begin{frame}{Exponentes: Ley de Potencias y Exponencial}
    \begin{columns}
      \begin{column}{0.5\textwidth}
        \centering
        \medium{Exponente de la Ley de Potencias}
        \vspace{0.2cm}
        \begin{figure}
          \centering
          \includegraphics[width=\linewidth]{imagenv25/alpha_vs_phi.png}
          \label{fig:alpha_vs_phi}
        \end{figure}  
      \end{column}
      \begin{column}{0.5\textwidth}
        \centering
        \medium{Exponente de Exponencial}
        \vspace{0.2cm}
        \begin{figure}
          \centering
          \includegraphics[width=\linewidth]{imagenv25/lambda_vs_phi.png}
          \label{fig:lambda_vs_phi}
        \end{figure}
      \end{column}
    \end{columns}
  \end{frame}
  
  % ------------------ Resultados L = 8 ------------------

  \subsection{Resultados L 8}
  \begin{frame}
    \centering
    \colorbox{green!20!gray!30}{%
      \begin{minipage}{0.85\textwidth}
        \centering
        \vspace{0.5cm}
        {\Large\ Resultados L 8}
        \vspace{0.5cm}
      \end{minipage}
    }
  \end{frame}

  \subsubsection{Animaciones}
  \begin{frame}{Animaciones}
  \end{frame}
  \begin{frame}{Animaciones}
    \begin{columns}
      \begin{column}{0.5\textwidth}
        \medium {N = 72}
        \includegraphics[width=1\linewidth]{imagenl8/ani_N189_L8.png}
        \url {https://www.youtube.com/watch?v=JMb-PDV1A3w}
      \end{column}
      \begin{column}{0.5\textwidth}
        \medium {N = 220}
        \includegraphics[width=1\linewidth]{imagenl8/ani_N389_L8.png}
        \url {https://www.youtube.com/watch?v=1X6egGP7IyQ}
      \end{column}
    \end{columns}
  \end{frame}

  \subsubsection{Contactos Únicos con Agente Fijo}
  \begin{frame}{Contactos Únicos con Agente Fijo}
    \begin{figure}
      \centering
      \includegraphics[width=1\linewidth]{imagenl8/accumulated_contacts_N_9_189_389.png}
      \label{fig:placeholder}
    \end{figure}
    \small {para \(L=8.0m\), $N \in $[10, 390]}
  \end{frame}

  \subsubsection{Scanning Rate}
  \begin{frame}{Scanning Rate}
    \begin{figure}
      \centering
      \includegraphics[width=1\linewidth]{imagenl8/scanning_rate_vs_phi.png}
      \label{fig:placeholder}
    \end{figure}
  \end{frame}

  \subsubsection{Aproximaciones}

  \begin{frame}{Aproximaciones}
    \begin{columns}
      \begin{column}{0.5\textwidth}
        \medium {Aproximación Power Law}
        \begin{figure}
          \centering
          \small{N=119}
          \includegraphics[width=1\linewidth]{imagenl8/distribution_N119.png}
          \label{fig:placeholder}
        \end{figure}
      \end{column}
      \begin{column}{0.5\textwidth}
        \medium {Aproximación Exponencial}
        \begin{figure}
          \centering
          \small{N=299}
          \includegraphics[width=1\linewidth]{imagenl8/distribution_N299.png}
          \label{fig:placeholder}
        \end{figure}
      \end{column}
    \end{columns}
  \end{frame}

  \subsubsection{Exponentes Aproximaciones}
  \begin{frame}{Exponentes}
    \begin{columns}
      \begin{column}{0.5\textwidth}
        \centering
        \medium{Exponente de la Ley de Potencias}
        \vspace{0.2cm}
        \begin{figure}
          \centering
          \includegraphics[width=\linewidth]{imagenl8/alpha_vs_phi.png}
          \label{fig:alpha_vs_phi}
        \end{figure}  
      \end{column}
      \begin{column}{0.5\textwidth}
        \centering
        \medium{Exponente de Exponencial}
        \vspace{0.2cm}
        \begin{figure}
          \centering
          \includegraphics[width=\linewidth]{imagenl8/lambda_vs_phi.png}
          \label{fig:lambda_vs_phi}
        \end{figure}
      \end{column}
    \end{columns}
  \end{frame}


  % ------------------ Conclusiones ------------------
  \section{Conclusiones}
  \begin{frame}
    \centering
    \colorbox{green!20!gray!30}{%
      \begin{minipage}{0.85\textwidth}
        \centering
        \vspace{0.5cm}
        {\Large\ Conclusiones}
        \vspace{0.5cm}
      \end{minipage}
    }
  \end{frame}

  \subsection{Conclusiones}
  \begin{frame}{Conclusiones}
    \Large {Estudio 1:}
    \begin{itemize}
      \item $ \phi \sim 0.25 $ es el valor óptimo para obtener la mayor cantidad de contactos.
      \vspace{0.3cm}
      \item El Scanning Rate está en su mayor punto para densidades intermedias.
      \vspace{0.3cm}
      \item Una mayor velocidad deseada no afecta significativamente el Scanning Rate.
      \vspace{0.3cm}
      \item Un mayor tamaño del sistema manteniendo r constante, desplaza el pico hacia una mayor densidad.
    \end{itemize}
  \end{frame}
  \begin{frame}{Conclusiones}
    \Large {Estudio 2:}
    \begin{itemize}
      \item Algunos casos presentan un ajuste al modelo de Ley de Potencias.
      \vspace{0.3cm}
      \item La mayoría de las configuraciones se ajustan mejor a una distribucion Exponencial.
      \vspace{0.3cm}
      \item Las mismas conclusiones se obtienen variando la velocidad deseada y el tamaño del sistema.
    \end{itemize}
  \end{frame}
  % \subsection{Conclusiones V 1.7 y L 8.0}
  % \begin{frame}{Conclusiones}
  %   \Large {Estudio 1:}
  %   \begin{itemize}
  %     \item El Scanning Rate aumenta con la densidad de peatones.
  %     \vspace{0.3cm}
  %     \item El ajuste por Ley de Potencias es mejor para bajas densidades, mientras que
  %   \end{itemize}
  %   \Large {Estudio 2:}
  %   \begin{itemize}
  %     \item El aumento de la velocidad deseada incrementa el Scanning Rate.
  %     \vspace{0.3cm}
  %     \item El ajuste por Ley de Potencias es mejor para bajas densidades, mientras que
  %   \end{itemize}
  % \end{frame}
  % \subsection{Conclusiones V 2.5 y L 6.0}
  % \begin{frame}{Conclusiones}
  %   \Large {Estudio 1:}
  %   \begin{itemize}
  %     \item El Scanning Rate aumenta con la densidad de peatones.
  %     \vspace{0.3cm}
  %     \item El ajuste por Ley de Potencias es mejor para bajas densidades, mientras que
  %   \end{itemize}
  %   \Large {Estudio 2:}
  %   \begin{itemize}
  %     \item El aumento de la velocidad deseada incrementa el Scanning Rate.
  %     \vspace{0.3cm}
  %     \item El ajuste por Ley de Potencias es mejor para bajas densidades, mientras que
  %   \end{itemize}
  % \end{frame}



% Cierre
  \begin{frame}
    \centering
    \Huge ¡Gracias por su atención!
  \end{frame}

\end{document}